\section{Fokker--Planck 方程描述密度的流动}
\label{sec:06-Random-Processes/06-Diffusion-and-SDE/02}

跑一次扩散模拟，你得到一条谁也预测不了的轨迹；换个随机种子再跑，轨迹面目全非。可要是跑上一万次，把终点画成直方图，这张直方图每次都几乎一模一样——换种子、换机器、换实现语言都不改变它的形状。

单条路径不可复现，路径的\emph{分布}高度可复现。上一节把注意力全押在一条轨迹上，而那条轨迹处处不可微、看不出规律。本节把镜头拉远，只看密度 \(p(x,t)\)。奖赏很直接：密度不但光滑，还服从一个普通的偏微分方程，里面没有一点随机性。

这不是新招式。离散 Markov 链早就这么干过：单条链的跳法随机，行分布却老老实实按 \(\mu_n=\mu_0P^n\) 演化。状态从离散换成连续之后，矩阵让位给微分算子，思想照搬。

\subsection{先读懂一小步长什么样}

形式地写下噪声驱动的演化

\[
dX_t=\mu(X_t)\,dt+\sigma(X_t)\,dW_t,
\]

它的严格含义要等下一节，本节只用局部读法：从位置 \(x\) 出发，经过很短的 \(\Delta t\)，

\[
\Delta X\approx\mu(x)\,\Delta t+\sigma(x)\sqrt{\Delta t}\,Z,
\qquad Z\sim\mathcal N(0,1).
\]

两项的量级不同，这是上一节那个 \(\sqrt{\Delta t}\) 的直接延续：均值被 \(\mu\) 挪动 \(O(\Delta t)\)，散布被 \(\sigma\) 撑开 \(O(\sqrt{\Delta t})\)。\(\mu\) 叫漂移系数，\(\sigma\) 叫扩散系数。和上一节唯一的区别是：方向可以有偏好，步子大小可以随位置变。

\subsection{质量守恒把方程逼出来}

概率质量不生不灭，某点密度的变化只能来自净流入。记 \(J(x,t)\) 为概率通量——单位时间内向右净通过 \(x\) 的概率质量，于是

\[
\frac{\partial p}{\partial t}=-\frac{\partial J}{\partial x}.
\]

这是一维连续性方程，和流体力学里的质量守恒一模一样。剩下的问题只有一个：\(J\) 由什么组成。

漂移那部分像河水载着染料，整团密度以速度 \(\mu(x)\) 平移，搬运的是密度本身，\(J_{\mathrm{drift}}=\mu(x)\,p\)。

扩散那部分要绕一下。单个粒子向左向右没有偏好，群体却有净效果：密度高的一侧走出来的粒子多，密度低的一侧走进去的少，净迁移与密度梯度反向，大小正比于局部散布强度。这就是 Fick 定律的直觉，\(\sigma\) 为常数时

\[
J_{\mathrm{diff}}=-\frac{\sigma^2}{2}\frac{\partial p}{\partial x}.
\]

若 \(\sigma\) 随位置变化，散布强度本身在变，正确的写法是 \(-\tfrac12\partial_x\big(\sigma(x)^2p\big)\)——导数要作用在整个 \(\sigma^2 p\) 上。这一步用物理直觉不容易分辨，严格来源是下一节的 Itô 公式：对任意光滑测试函数 \(\varphi\) 展开 \(\E[\varphi(X_{t+\Delta t})]\) 到二阶，再把导数分部积分搬到 \(p\) 上，就得到唯一正确的形式。合并两项代入连续性方程：

\[
\frac{\partial p}{\partial t}
=-\frac{\partial}{\partial x}\big(\mu(x)\,p\big)
+\frac12\frac{\partial^2}{\partial x^2}\big(\sigma(x)^2p\big).
\]

这就是 Fokker--Planck 方程，也叫 Kolmogorov 前向方程。给定初始密度 \(p(x,0)\)（确定性起点对应集中于一点的点质量），它把全部演化定死。

\begin{center}
\begin{tikzpicture}[x=1cm,y=1cm,font=\small,>=stealth]
  \draw[->,black!55] (-0.25,0) -- (4.15,0) node[right,font=\footnotesize] {\(t\)};
  \draw[->,black!55] (-0.25,-1.6) -- (-0.25,1.75) node[above,font=\footnotesize] {\(x\)};
  \draw[blue!35] (0.00,0.00) -- (0.10,-0.23) -- (0.20,0.00) -- (0.30,0.23) -- (0.40,0.23) -- (0.50,0.23) -- (0.60,0.46) -- (0.70,0.23) --
    (0.80,0.23) -- (0.90,0.46) -- (1.00,0.46) -- (1.10,0.23) -- (1.20,0.23) -- (1.30,0.23) -- (1.40,0.46) -- (1.50,0.46) --
    (1.60,0.23) -- (1.70,0.23) -- (1.80,0.23) -- (1.90,0.23) -- (2.00,0.23) -- (2.10,0.23) -- (2.20,0.23) -- (2.30,0.23) --
    (2.40,0.00) -- (2.50,-0.23) -- (2.60,-0.46) -- (2.70,-0.46) -- (2.80,-0.69) -- (2.90,-0.92) -- (3.00,-1.15) -- (3.10,-1.15) --
    (3.20,-1.15) -- (3.30,-1.38) -- (3.40,-1.38);
  \draw[blue!35] (0.00,0.00) -- (0.10,0.23) -- (0.20,0.46) -- (0.30,0.46) -- (0.40,0.46) -- (0.50,0.46) -- (0.60,0.46) -- (0.70,0.46) --
    (0.80,0.23) -- (0.90,0.23) -- (1.00,0.46) -- (1.10,0.69) -- (1.20,0.69) -- (1.30,0.69) -- (1.40,0.69) -- (1.50,0.69) --
    (1.60,0.69) -- (1.70,0.46) -- (1.80,0.69) -- (1.90,0.92) -- (2.00,0.69) -- (2.10,0.69) -- (2.20,0.69) -- (2.30,0.69) --
    (2.40,0.46) -- (2.50,0.46) -- (2.60,0.23) -- (2.70,0.23) -- (2.80,0.00) -- (2.90,-0.23) -- (3.00,-0.23) -- (3.10,-0.46) --
    (3.20,-0.23) -- (3.30,-0.23) -- (3.40,-0.46);
  \draw[blue!35] (0.00,0.00) -- (0.10,0.00) -- (0.20,0.23) -- (0.30,0.00) -- (0.40,0.00) -- (0.50,0.00) -- (0.60,-0.23) -- (0.70,-0.23) --
    (0.80,-0.23) -- (0.90,-0.46) -- (1.00,-0.46) -- (1.10,-0.46) -- (1.20,-0.46) -- (1.30,-0.23) -- (1.40,-0.23) -- (1.50,0.00) --
    (1.60,0.23) -- (1.70,0.00) -- (1.80,0.00) -- (1.90,0.00) -- (2.00,0.00) -- (2.10,0.00) -- (2.20,0.00) -- (2.30,-0.23) --
    (2.40,-0.23) -- (2.50,-0.23) -- (2.60,-0.23) -- (2.70,-0.23) -- (2.80,-0.23) -- (2.90,-0.46) -- (3.00,-0.46) -- (3.10,-0.69) --
    (3.20,-0.69) -- (3.30,-0.69) -- (3.40,-0.69);
  \draw[blue!70,thick] (0.00,0.00) -- (0.10,0.00) -- (0.20,0.00) -- (0.30,-0.23) -- (0.40,0.00) -- (0.50,0.00) -- (0.60,0.00) -- (0.70,0.00) --
    (0.80,-0.23) -- (0.90,0.00) -- (1.00,0.00) -- (1.10,0.00) -- (1.20,0.23) -- (1.30,0.00) -- (1.40,0.23) -- (1.50,0.23) --
    (1.60,0.23) -- (1.70,0.46) -- (1.80,0.46) -- (1.90,0.46) -- (2.00,0.69) -- (2.10,0.69) -- (2.20,0.46) -- (2.30,0.46) --
    (2.40,0.69) -- (2.50,0.69) -- (2.60,0.46) -- (2.70,0.46) -- (2.80,0.46) -- (2.90,0.46) -- (3.00,0.69) -- (3.10,0.69) --
    (3.20,0.69) -- (3.30,0.46) -- (3.40,0.23);
  \draw[black!30,dashed] (0.55,-1.55) -- (0.55,1.55);
  \draw[black!30,dashed] (1.70,-1.55) -- (1.70,1.55);
  \draw[black!30,dashed] (3.40,-1.55) -- (3.40,1.55);
  \draw[red!70,thick] plot[domain=-1.55:1.55,samples=60] ({0.55+0.90*exp(-\x*\x/0.292)},{\x});
  \draw[red!70,thick] plot[domain=-1.55:1.55,samples=60] ({1.70+0.51*exp(-\x*\x/0.902)},{\x});
  \draw[red!70,thick] plot[domain=-1.55:1.55,samples=60] ({3.40+0.36*exp(-\x*\x/1.805)},{\x});
  \fill[red!80] (0,0) circle (2pt);
  \node[left,font=\footnotesize,red!80] at (-0.05,0.28) {点质量};
  \node[below,font=\footnotesize] at (0.55,-1.62) {\(t_1\)};
  \node[below,font=\footnotesize] at (1.70,-1.62) {\(t_2\)};
  \node[below,font=\footnotesize] at (3.40,-1.62) {\(t_3\)};
\end{tikzpicture}

\emph{四条蓝色轨迹来自同一个起点、同一个方程，彼此毫无相似之处；粗线那条也没有任何特殊地位。红色曲线是三个时刻上的密度截面，由尖峰逐步塌成矮胖的鼓包，宽度按 \(\sqrt{t}\) 长大而峰高按 \(1/\sqrt t\) 落下，面积始终是 \(1\)。图里唯一确定、唯一可复现的东西就是这三条红线，而 Fokker--Planck 方程算的正是它们。}
\end{center}

\subsection{它和离散链的前向方程是同一件事}

\hyperref[sec:06-Random-Processes/05-Continuous-Time-Markov-Chains-and-Queues/02]{``Kolmogorov 方程与矩阵指数''}里，行分布满足主方程 \(\mu'(t)=\mu(t)Q\)，写成分量就是固定终点 \(j\)，按``从哪里跳进来''汇总流入。Fokker--Planck 是同一个前向视角，只不过固定的是位置 \(x\)，问密度被周围的漂移和扩散怎样送进来。

{\def\LTcaptype{none}
\begin{longtable}{@{}lll@{}}
\toprule\noalign{}
 & 离散状态（CTMC） & 连续状态（扩散） \\
\midrule\noalign{}
\endhead
\bottomrule\noalign{}
\endlastfoot
演化对象 & 行向量 \(\mu(t)\) & 密度 \(p(\cdot,t)\) \\
局部规则 & 速率矩阵 \(Q\) & 漂移 \(\mu(x)\)、扩散 \(\sigma(x)\) \\
前向方程 & \(\mu'(t)=\mu(t)Q\) & 上面那个偏微分方程 \\
闭式解 & \(\mu(0)e^{tQ}\) & 一般没有，靠数值解 \\
\end{longtable}
}

求和换成了微分，矩阵指数的位置被微分算子生成的演化占据。限制写在最后一行：能写下方程不等于能解出方程——这是离开有限状态之后遇到的第一道坎。

\subsection{纯扩散只会摊平}

取 \(\mu=0\)、\(\sigma\) 为常数，方程退化成热方程 \(\partial_tp=\tfrac{\sigma^2}{2}\partial_x^2p\)。从原点的点质量出发，解是 Gaussian 热核

\[
p(x,t)=\frac{1}{\sqrt{2\pi\sigma^2t}}\exp\!\left(-\frac{x^2}{2\sigma^2t}\right),
\]

也就是 \(\mathcal N(0,\sigma^2t)\)——恰是上一节布朗运动在时刻 \(t\) 的边缘分布。路径极限和偏微分方程描述的是同一个对象，一个看单条轨迹，一个看群体密度。

方差按 \(t\) 线性长大，密度被无限摊薄，因此纯扩散在整条直线上没有平稳分布。摊平这件事一旦没人约束，质量就一路散到无穷远。要有平稳形态，必须有一股力气把它拽回来。

\subsection{摊平与回复的平衡}

给它加一个指向原点的线性漂移：

\[
dX_t=-\theta X_t\,dt+\sigma\,dW_t,\qquad\theta>0,
\]

这是 Ornstein--Uhlenbeck（OU）过程，\(\abs{X_t}\) 越大回拉越狠，像根弹簧。找平稳密度不必解偏微分方程，令 \(\partial_tp=0\) 并取``通量为零''的自然边界（概率不许持续流向无穷远）：

\[
J=-\theta x\,p-\frac{\sigma^2}{2}\frac{dp}{dx}=0
\;\Longrightarrow\;
\frac{d}{dx}\log p=-\frac{2\theta}{\sigma^2}x,
\]

积分得 \(p(x)\propto\exp\big(-\theta x^2/\sigma^2\big)\)，即 \(\mathcal N\big(0,\sigma^2/(2\theta)\big)\)。

扩散往外摊，回复往里拽，两股力量在某个宽度上打平。哪怕 \(\theta\) 很小也够用，只是平衡宽度 \(\sigma^2/(2\theta)\) 变得很大。这个``扩散加回复''的配方是制造可控平稳分布的基本手法，后面两节会看到扩散生成模型正是靠它把任意数据分布洗成一个已知的 Gaussian。

\subsection{把想要的分布写进漂移}

上面那步倒推值得再走一遍，因为它藏着本单元最关键的一次身份转换。零通量条件说的是

\[
\mu(x)=\frac{\sigma^2}{2}\frac{d}{dx}\log p(x),
\]

漂移必须正比于\emph{平稳密度的对数导数}。反过来读：给定你想采样的目标密度 \(\pi\)，只要把漂移设成 \(\nabla\log\pi\)，扩散设成对应强度，这个 SDE 的平稳分布就是 \(\pi\)。写成常见形式，

\[
dX_t=\nabla\log\pi(X_t)\,dt+\sqrt{2}\,dW_t,
\]

叫 Langevin 动力学。跑得够久，轨迹上的点就是 \(\pi\) 的样本。

这里有两件事值得停一下。第一，\(\nabla\log\pi\) 这个量马上会有名字——得分（score），它是后面两节的主角。此刻它以最朴素的身份登场：\emph{让扩散停在指定分布上所需要的那股漂移}。第二，\(\pi\) 只需知道到差一个常数倍：若 \(\pi=e^{-U}/Z\)，那么 \(\nabla\log\pi=-\nabla U\)，归一化常数 \(Z\) 求导时消失得无影无踪。配分函数 \(Z\) 通常是概率建模里最难算的东西（见\hyperref[sec:13-Machine-Learning/08-Sampling-and-Inference/06]{``配分函数把能量变成全局概率''}），而 Langevin 采样从头到尾不需要它。

把能量写活一点，加上温度参数：

\[
dX_t=-\nabla U(X_t)\,dt+\sqrt{2T}\,dW_t
\quad\Longrightarrow\quad
\pi_T(x)\propto e^{-U(x)/T}.
\]

\begin{center}
\begin{tikzpicture}[x=1.7cm,y=1cm,font=\small,>=stealth]
  \draw[black!45] (-1.85,0) -- (1.85,0);
  \draw[black!55,thick] plot[domain=-1.62:1.62,samples=70] (\x,{1.75+0.35*(\x*\x-1)*(\x*\x-1)});
  \node[right,font=\footnotesize,black!70] at (1.5,3.0) {\(U(x)\)};
  \fill[black!45] (-1,1.75) circle (1.4pt);
  \fill[black!45] (1,1.75) circle (1.4pt);
  \draw[blue!75,thick] plot[domain=-1.62:1.62,samples=90] (\x,{1.20*exp(-(\x*\x-1)*(\x*\x-1)/0.15)});
  \draw[red!70,thick] plot[domain=-1.75:1.75,samples=70] (\x,{0.70*exp(-(\x*\x-1)*(\x*\x-1)/1.0)});
  \node[right,font=\footnotesize,blue!75] at (1.05,1.15) {低温 \(T=0.15\)};
  \node[right,font=\footnotesize,red!70] at (1.25,0.72) {高温 \(T=1\)};
  \node[below,font=\footnotesize] at (0,-0.25) {\(x\)};
\end{tikzpicture}

\emph{同一个双阱能量 \(U\)，温度不同给出完全不同的平稳密度。高温时密度几乎摊成一片，粒子轻松翻越中间的势垒，采样跑得动却不精细；低温时密度锁死在两个谷底，样本很准，可一旦掉进左边的坑就几乎出不来。退火采样的做法是先高温后低温：让粒子在高温下找到该去哪个区域，再降温把它压到谷底。下一节的``噪声水平从大到小''和这张图是同一件事。}
\end{center}

温度调的是探索与精确之间的天平，这跟\hyperref[sec:13-Machine-Learning/08-Sampling-and-Inference/02]{``MCMC 在无法独立采样时构造平稳链''}里步长与接受率的权衡是同一类困难：链混合得快就采得糙，采得细就容易困在局部。区别在于 Langevin 用的是梯度信息，它知道往哪边走密度更高，因而在高维里比盲目游走的 Metropolis 有效得多。

\subsection{边界}

Fokker--Planck 方程假设过程真的是连续扩散，路径不许跳。有跳跃时（重尾冲击、突发事件）右端要补一个积分项，那属于跳扩散的理论，本单元不涉及。系数太坏——不光滑、增长过快——时演化可能不存在，或者在有限时间里爆掉，下一节会给出常用的充分条件。还有一个容易忽视的前提：上面所有推导都默认密度存在。初值是点质量时 \(t=0\) 那一刻并没有密度，是扩散在任意正时刻立刻把它抹开的，这依赖 \(\sigma\) 不退化；\(\sigma\) 在某处为零，那里的密度可以一直是奇异的。

密度视角能回答``分布怎样演化''，回答不了``第一次碰到某个阈值是什么时候''这类只关于单条路径的问题，也没法处理路径泛函。更要紧的是，\(dX=\mu\,dt+\sigma\,dW\) 至今还只是一个记号。下一节把这个记号兑现：先定义 \(\int\cdot\,dW\)，再看链式法则被 \(\sqrt{\Delta t}\) 改成了什么样子。
